\chapter{Introducción}
\label{cap:introduccion}

\section{Motivación}

La epilepsia es una de las enfermedades neurológicas crónicas más frecuentes: afecta a unos 50 millones de personas en el mundo. Se define por la predisposición duradera a generar crisis epilépticas, episodios transitorios causados por una actividad neuronal anormalmente excesiva o sincrónica \parencite{fisher2014}. El electroencefalograma (EEG) de superficie es la herramienta central para el diagnóstico y el seguimiento, pero su revisión es manual y lenta: un registro continuo de 24 horas puede requerir horas de lectura de un especialista.

Un detector automático de crisis tiene que cumplir dos objetivos en conflicto. Debe ser \emph{sensible}, porque una crisis no detectada puede tener consecuencias graves, y \emph{específico}, porque un sistema con muchas falsas alarmas termina siendo ignorado. En la práctica clínica la especificidad suele expresarse como la cantidad de \emph{falsas alarmas por hora}.

\pendiente{ampliar a 2--3 páginas: impacto clínico, monitoreo prolongado, dispositivos portátiles, por qué el aprendizaje automático.}

\section{Planteamiento del problema}

Se aborda la detección de crisis como un problema de clasificación binaria de ventanas de EEG de 4 segundos. El problema tiene tres dificultades: (i) un fuerte desbalance de clases, porque las crisis ocupan menos del 1~\% del tiempo registrado; (ii) una gran variabilidad entre pacientes, porque la forma de una crisis cambia de una persona a otra; y (iii) la necesidad de \emph{generalizar a pacientes nuevos}, lo que obliga a evaluar con pacientes que el modelo no vio durante el entrenamiento.

\section{Objetivos}

\paragraph{Objetivo general.} Construir y evaluar un detector de crisis epilépticas basado en características espectrales del EEG y redes neuronales, y estudiar de forma controlada cómo influyen el optimizador y la función de activación en el entrenamiento y en el rendimiento.

\paragraph{Objetivos específicos.}
\begin{enumerate}
  \item Modelar matemáticamente la señal EEG y fundamentar, con la transformada de Fourier, las características espectrales que distinguen una crisis.
  \item Demostrar que la entropía espectral alcanza su máximo global en la distribución uniforme del simplex y verificar empíricamente que disminuye durante las crisis.
  \item Comparar seis optimizadores en un problema convexo y en uno no convexo, incluyendo el método del simplex de Nelder-Mead como referencia sin derivadas.
  \item Comparar cinco funciones de activación y su interacción con el optimizador y con la profundidad de la red.
  \item Evaluar el detector con curva ROC, AUC, matriz de confusión, sensibilidad y especificidad, eligiendo el umbral como un problema de optimización con restricciones.
  \item Controlar el sobreajuste con partición por paciente, regularización y \emph{early stopping}.
\end{enumerate}

\section{Hipótesis}
\label{sec:hipotesis}

\begin{hipotesis}\label{h:convexo}
En la regresión logística regularizada, que es un problema estrictamente convexo, todos los optimizadores convergen al mismo óptimo global y solo difieren en la velocidad.
\end{hipotesis}
\begin{hipotesis}\label{h:noconvexo}
En el perceptrón multicapa, que es un problema no convexo, la variabilidad del resultado entre inicializaciones es significativamente mayor que en el caso convexo.
\end{hipotesis}
\begin{hipotesis}\label{h:profundidad}
Al aumentar la profundidad, las redes con activaciones saturantes (sigmoide y tanh) se degradan más que las de la familia ReLU y GELU, por desvanecimiento del gradiente.
\end{hipotesis}
\begin{hipotesis}\label{h:adaptativos}
Los optimizadores adaptativos (RMSProp, Adam) son menos sensibles a la elección de la función de activación que SGD.
\end{hipotesis}
\begin{hipotesis}\label{h:entropia}
La entropía espectral del EEG es significativamente menor durante las crisis que fuera de ellas.
\end{hipotesis}

\section{Organización del trabajo}

El capítulo~\ref{cap:senal} presenta la fisiología del EEG y un modelo matemático de la señal. El capítulo~\ref{cap:fourier} desarrolla el análisis de Fourier y las características espectrales, incluida la entropía espectral en el simplex. El capítulo~\ref{cap:aprendizaje} formaliza el problema de aprendizaje y el sobreajuste. Los capítulos~\ref{cap:modelos} y~\ref{cap:optimizacion} tratan los modelos, las funciones de activación y los métodos de optimización. El capítulo~\ref{cap:evaluacion} define las métricas y la elección del umbral. El capítulo~\ref{cap:metodologia} describe el diseño experimental, el capítulo~\ref{cap:resultados} presenta los resultados y el capítulo~\ref{cap:conclusiones} los discute.

Todo el código está en el notebook \texttt{TP\_Final\_Crisis\_EEG.ipynb} y en los módulos de la carpeta \texttt{src/}, y puede ejecutarse localmente o en Google Colab.
